\documentclass[border=5pt]{standalone}
\usepackage[T1]{fontenc}
\usepackage{tikz}
\usetikzlibrary{arrows.meta,positioning,intersections,calc}

\definecolor{garnet}{HTML}{73000A}
\definecolor{atlantic}{HTML}{466A9F}
\definecolor{horseshoe}{HTML}{65780B}
\definecolor{bk70}{HTML}{5C5C5C}
\definecolor{bk30}{HTML}{C7C7C7}
\definecolor{warmgrey}{HTML}{676156}

\begin{document}
\begin{tikzpicture}[scale=1.0]

  \pgfmathsetmacro{\beamangle}{\PARAM}  % tikzgif sweeps this
  \def\beamlen{5.5}

  % Bounding box — adjust these to control size and centering
  \useasboundingbox (-1.5,-3) rectangle (6.5,3);
  \draw[horseshoe, thick] (-1.5,-3) rectangle (6.5,3);

  % Step 1: Draw objects — spread around the LiDAR, all within box bounds
  \draw[thick, warmgrey, fill=warmgrey!25] (3.5,1.5) circle (0.4);          % right-upper
  \draw[thick, warmgrey, fill=warmgrey!25] (4.5,-0.5) +(-0.7,-0.3) rectangle +(0.7,0.3); % right-mid
  \draw[thick, warmgrey, fill=warmgrey!25]
    (2.5,-1.8) +(-0.4,0) -- +(0,0.6) -- +(0.4,0) -- cycle;                % below-right
  \draw[thick, warmgrey, fill=warmgrey!25] (5.0,2.2) circle (0.3);          % far right-upper
  \draw[thick, warmgrey, fill=warmgrey!25] (3.0,-2.3) +(-0.5,-0.2) rectangle +(0.5,0.2); % below
  \draw[thick, warmgrey, fill=warmgrey!25]
    (-0.5,2.0) +(-0.3,0) -- +(0,0.5) -- +(0.3,0) -- cycle;                % above LiDAR
  \draw[thick, warmgrey, fill=warmgrey!25] (1.5,-2.0) circle (0.35);        % below-left

  % Step 2: Invisible paths for intersection detection
  % Beam extends far enough to always hit the box wall
  \path[name path=beam] (0,0) -- (\beamangle:10);
  \path[name path=box] (-1.5,-3) rectangle (6.5,3);
  \path[name path=circ] (3.5,1.5) circle (0.4);
  \path[name path=rect] (4.5,-0.5) +(-0.7,-0.3) rectangle +(0.7,0.3);
  \path[name path=tri]
    (2.5,-1.8) +(-0.4,0) -- +(0,0.6) -- +(0.4,0) -- cycle;
  \path[name path=circ2] (5.0,2.2) circle (0.3);
  \path[name path=rect2] (3.0,-2.3) +(-0.5,-0.2) rectangle +(0.5,0.2);
  \path[name path=tri2]
    (-0.5,2.0) +(-0.3,0) -- +(0,0.5) -- +(0.3,0) -- cycle;
  \path[name path=circ3] (1.5,-2.0) circle (0.35);

  % Step 3: Find closest intersection across all objects
  \def\mindist{999}
  \def\hitx{0}
  \def\hity{0}
  \def\gothit{0}

  % Check circle
  \path[name intersections={of=beam and circ, total=\t}]
    \pgfextra{
      \ifnum\t>0
        \foreach \i in {1,...,\t} {
          \pgfpointanchor{intersection-\i}{center}
          \pgfgetlastxy{\ix}{\iy}
          \pgfmathparse{veclen(\ix,\iy)}
          \let\d\pgfmathresult
          \pgfmathparse{\d < \mindist ? 1 : 0}
          \ifdim\pgfmathresult pt>0.5pt
            \global\edef\mindist{\d}
            \global\edef\hitx{\ix}
            \global\edef\hity{\iy}
            \global\def\gothit{1}
          \fi
        }
      \fi
    };

  % Check rectangle
  \path[name intersections={of=beam and rect, total=\t}]
    \pgfextra{
      \ifnum\t>0
        \foreach \i in {1,...,\t} {
          \pgfpointanchor{intersection-\i}{center}
          \pgfgetlastxy{\ix}{\iy}
          \pgfmathparse{veclen(\ix,\iy)}
          \let\d\pgfmathresult
          \pgfmathparse{\d < \mindist ? 1 : 0}
          \ifdim\pgfmathresult pt>0.5pt
            \global\edef\mindist{\d}
            \global\edef\hitx{\ix}
            \global\edef\hity{\iy}
            \global\def\gothit{1}
          \fi
        }
      \fi
    };

  % Check triangle
  \path[name intersections={of=beam and tri, total=\t}]
    \pgfextra{
      \ifnum\t>0
        \foreach \i in {1,...,\t} {
          \pgfpointanchor{intersection-\i}{center}
          \pgfgetlastxy{\ix}{\iy}
          \pgfmathparse{veclen(\ix,\iy)}
          \let\d\pgfmathresult
          \pgfmathparse{\d < \mindist ? 1 : 0}
          \ifdim\pgfmathresult pt>0.5pt
            \global\edef\mindist{\d}
            \global\edef\hitx{\ix}
            \global\edef\hity{\iy}
            \global\def\gothit{1}
          \fi
        }
      \fi
    };

  % Check circle2
  \path[name intersections={of=beam and circ2, total=\t}]
    \pgfextra{
      \ifnum\t>0
        \foreach \i in {1,...,\t} {
          \pgfpointanchor{intersection-\i}{center}
          \pgfgetlastxy{\ix}{\iy}
          \pgfmathparse{veclen(\ix,\iy)}
          \let\d\pgfmathresult
          \pgfmathparse{\d < \mindist ? 1 : 0}
          \ifdim\pgfmathresult pt>0.5pt
            \global\edef\mindist{\d}
            \global\edef\hitx{\ix}
            \global\edef\hity{\iy}
            \global\def\gothit{1}
          \fi
        }
      \fi
    };

  % Check rect2
  \path[name intersections={of=beam and rect2, total=\t}]
    \pgfextra{
      \ifnum\t>0
        \foreach \i in {1,...,\t} {
          \pgfpointanchor{intersection-\i}{center}
          \pgfgetlastxy{\ix}{\iy}
          \pgfmathparse{veclen(\ix,\iy)}
          \let\d\pgfmathresult
          \pgfmathparse{\d < \mindist ? 1 : 0}
          \ifdim\pgfmathresult pt>0.5pt
            \global\edef\mindist{\d}
            \global\edef\hitx{\ix}
            \global\edef\hity{\iy}
            \global\def\gothit{1}
          \fi
        }
      \fi
    };

  % Check tri2
  \path[name intersections={of=beam and tri2, total=\t}]
    \pgfextra{
      \ifnum\t>0
        \foreach \i in {1,...,\t} {
          \pgfpointanchor{intersection-\i}{center}
          \pgfgetlastxy{\ix}{\iy}
          \pgfmathparse{veclen(\ix,\iy)}
          \let\d\pgfmathresult
          \pgfmathparse{\d < \mindist ? 1 : 0}
          \ifdim\pgfmathresult pt>0.5pt
            \global\edef\mindist{\d}
            \global\edef\hitx{\ix}
            \global\edef\hity{\iy}
            \global\def\gothit{1}
          \fi
        }
      \fi
    };

  % Check circ3
  \path[name intersections={of=beam and circ3, total=\t}]
    \pgfextra{
      \ifnum\t>0
        \foreach \i in {1,...,\t} {
          \pgfpointanchor{intersection-\i}{center}
          \pgfgetlastxy{\ix}{\iy}
          \pgfmathparse{veclen(\ix,\iy)}
          \let\d\pgfmathresult
          \pgfmathparse{\d < \mindist ? 1 : 0}
          \ifdim\pgfmathresult pt>0.5pt
            \global\edef\mindist{\d}
            \global\edef\hitx{\ix}
            \global\edef\hity{\iy}
            \global\def\gothit{1}
          \fi
        }
      \fi
    };

  % Find where beam hits the bounding box wall
  \def\wallx{0}
  \def\wally{0}
  \path[name intersections={of=beam and box, total=\tw}]
    \pgfextra{
      \def\walldist{999}
      \ifnum\tw>0
        \foreach \i in {1,...,\tw} {
          \pgfpointanchor{intersection-\i}{center}
          \pgfgetlastxy{\ix}{\iy}
          \pgfmathparse{veclen(\ix,\iy)}
          \let\d\pgfmathresult
          \pgfmathparse{\d < \walldist ? 1 : 0}
          \ifdim\pgfmathresult pt>0.5pt
            \global\edef\walldist{\d}
            \global\edef\wallx{\ix}
            \global\edef\wally{\iy}
          \fi
        }
      \fi
    };

  % Step 4: Draw beam
  % Hit object -> garnet line from L to object, blue dot on object
  % Miss all  -> green line from L to wall, blue dot on wall
  \ifnum\gothit=1
    \draw[thick, garnet] (0,0) -- (\hitx, \hity);
    \fill[atlantic] (\hitx, \hity) circle (0.1);
  \else
    \draw[thick, horseshoe] (0,0) -- (\wallx, \wally);
    \fill[atlantic] (\wallx, \wally) circle (0.1);
  \fi

  % Step 5: LiDAR unit (on top)
  \draw[very thick, atlantic, fill=white] (0,0) circle (0.5);
  \node[font=\sffamily\bfseries, color=atlantic] at (0,0) {L};

\end{tikzpicture}
\end{document}
